\documentclass[aspectratio=169, 11pt, xcolor={table}]{beamer}
% Shared style for the course slides, based on presentation_text/main.tex.
% Dark blue removed: every slide (cover included) uses the light background
% of the template's content slides, with teal + orange as accents.

\usepackage{amsmath}
\usepackage{amssymb}
\usepackage{booktabs}
\usepackage{colortbl}
\usepackage{tabularx}
\usepackage{tikz}
\usepackage[skins]{tcolorbox}
\usepackage{fontawesome5}
\usepackage{polyglossia}

\usetikzlibrary{positioning, calc, arrows.meta, shapes.geometric}

%------------------------------------------------
% Languages and fonts (compile from slides/ with XeLaTeX or tectonic)
%------------------------------------------------
\setdefaultlanguage[numerals=maghrib]{arabic}
\setotherlanguage{english}
\usefonttheme{serif}

\newfontfamily\arabicfont[
    Script=Arabic, Scale=0.98, Path=../fonts/, Extension=.ttf,
    UprightFont=*-Regular, BoldFont=*-Bold,
    ItalicFont=*-Italic, SlantedFont=*-Italic,
    BoldItalicFont=*-BoldItalic, BoldSlantedFont=*-BoldItalic]{Amiri}
\newfontfamily\arabicfontsf[
    Script=Arabic, Scale=0.98, Path=../fonts/, Extension=.ttf,
    UprightFont=*-Regular, BoldFont=*-Bold,
    ItalicFont=*-Italic, SlantedFont=*-Italic,
    BoldItalicFont=*-BoldItalic, BoldSlantedFont=*-BoldItalic]{Amiri}
% Heading font
\newfontfamily\khatt[Script=Arabic, Scale=1.05, Path=../fonts/,
    BoldFont=ae_AlBattar.ttf]{ae_AlBattar.ttf}

%------------------------------------------------
% Colors (template palette without the navy tones)
%------------------------------------------------
\definecolor{brandTeal}{HTML}{00B4D8}
\definecolor{brandOrange}{HTML}{FF7043}
\definecolor{brandLight}{HTML}{F8F9FA}
\definecolor{brandDark}{HTML}{263238}
\definecolor{brandMuted}{HTML}{6C757D}
\colorlet{tealDeep}{brandTeal!70!black}

\setbeamercolor{background canvas}{bg=brandLight}
\setbeamercolor{normal text}{fg=brandDark}
\setbeamercolor{structure}{fg=brandTeal}
\setbeamertemplate{navigation symbols}{}

% beamer.cls sets \raggedright globally; under RTL that centers short lines
% or pushes them left, so go back to justified text (flush right).
\rightskip=0pt\relax
% Lists call \raggedright too; make it a no-op for RTL slides.
\AtBeginDocument{\renewcommand{\raggedright}{\rightskip=0pt\relax}}

%------------------------------------------------
% Frame title, footer, lists
%------------------------------------------------
\setbeamertemplate{frametitle}{%
    \vspace{0.35cm}
    \noindent{\khatt\bfseries\Large\color{brandDark}\insertframetitle}
    \par
    \vskip-2pt
    \noindent{\color{brandTeal}\rule{\textwidth}{1.5pt}}
    \vspace{0.05cm}
}

% Footer: page number only
\setbeamertemplate{footline}{%
    \hbox to \paperwidth{\hfill
        {\color{brandMuted}\scriptsize\LR{\insertframenumber\,/\,\inserttotalframenumber}}%
        \hspace{0.6cm}}%
    \vspace{0.25cm}}

\setbeamertemplate{itemize item}{\color{brandTeal}\small\faChevronLeft}
\setbeamertemplate{itemize subitem}{\color{brandMuted}\small\faAngleLeft}
\setbeamertemplate{enumerate item}{\color{tealDeep}\bfseries\insertenumlabel)}

% Tables: every cell framed by a visible border
\newcolumntype{R}{>{\raggedleft\arraybackslash}X}
\newcolumntype{C}{>{\centering\arraybackslash}X}
\newcommand{\ras}[1]{\cellcolor{brandTeal!20}\textbf{#1}}
% Colored text inside a table cell (\leavevmode keeps it on the row baseline)
\newcommand{\cell}[2]{\leavevmode\textcolor{#1}{#2}}
\renewcommand{\arraystretch}{1.35}
\setlength{\arrayrulewidth}{0.6pt}
\arrayrulecolor{brandMuted!70}

%------------------------------------------------
% Boxes (the two boxes of the template's 2nd slide)
%------------------------------------------------
\tcbset{
    brandbox/.style={
        enhanced, arc=4pt, outer arc=4pt, boxrule=1.5pt,
        coltitle=white, coltext=brandDark, fonttitle=\khatt\bfseries,
        toptitle=0.8mm, bottomtitle=0.8mm,
        top=1.2mm, bottom=1.2mm, left=3mm, right=3mm,
        before skip=4pt, after skip=4pt}
}
% Teal box: definitions, objectives, summaries
\newtcolorbox{tbox}[2][]{brandbox, colback=brandTeal!10,
    colframe=brandTeal, title={#2}, #1}
% Orange box: warnings
\newtcolorbox{obox}[2][]{brandbox, colback=brandOrange!10,
    colframe=brandOrange, title={#2}, #1}
% Card (template's 3rd slide, in teal)
\newtcolorbox{card}[2][]{enhanced, colback=brandTeal!6, colframe=brandTeal,
    arc=6pt, boxrule=1.2pt, halign=center, halign title=center,
    coltitle=white, coltext=brandDark, fonttitle=\khatt\bfseries,
    top=1.5mm, bottom=1.5mm, left=2mm, right=2mm, title={#2}, #1}
% Prompt text written to the tool
\newtcolorbox{amr}[1][brandTeal]{enhanced, arc=1mm, boxrule=0pt,
    colback=white, overlay={\fill[#1] (frame.north east) rectangle ([xshift=-2.5pt]frame.south east);},
    left=6pt, right=8pt, top=2pt, bottom=2pt,
    before skip=3pt, after skip=3pt}

%------------------------------------------------
% TikZ (Arabic right-to-left inside nodes)
%------------------------------------------------
\newcommand{\nodecolor}{}
\tikzset{
    every node/.append style={
        execute at begin node={\nodecolor\everypar{\beginR}\beginR},
        execute at end node=\endR},
    lawn/.code={\def\nodecolor{\color{#1}}},
    khatwa/.style={draw=#1, fill=#1!12, very thick, rounded corners=5pt,
        minimum height=0.95cm, text width=1.9cm, align=center,
        font=\khatt\bfseries\small},
    khatwa/.default=brandTeal,
    sahm/.style={-{Stealth[length=2.5mm]}, thick, color=brandMuted},
    rang/.style={draw=#1, fill=#1!12},
    raqm/.style={circle, fill=#1, lawn=white, inner sep=1.5pt,
        font=\small\bfseries},
}

% Horizontal chain of steps, right to left:
% \flow[width]{{a}, {b\\c}, {d}}
\newcommand{\flow}[2][1.9cm]{%
    \begin{tikzpicture}[node distance=5mm]
        \foreach \t [count=\i] in {#2} {
            \ifnum\i=1
                \node[khatwa, text width=#1] (f\i) {\t};
            \else
                \pgfmathtruncatemacro\p{\i-1}
                \node[khatwa, text width=#1, left=of f\p] (f\i) {\t};
                \draw[sahm] (f\p) -- (f\i);
            \fi
        }
    \end{tikzpicture}}

%------------------------------------------------
% Cover (template's 1st slide on the light background)
%------------------------------------------------
\newcommand{\cover}[3]{%
\begin{frame}[plain]
    \begin{tikzpicture}[remember picture, overlay]
        \draw[brandTeal, line width=1.5pt, opacity=0.3]
            (current page.south east) ++(-1,1) circle (4cm);
        \draw[brandTeal, line width=0.8pt, opacity=0.25]
            (current page.north west) ++(2,2) circle (6cm);
        \draw[brandOrange, line width=0.5pt, opacity=0.25]
            (current page.north west) ++(2,2) circle (7cm);
        \node[anchor=north east] at ($(current page.north east)+(-1.0,-0.8)$) {
            \begin{tikzpicture}[scale=0.35]
                \fill[brandTeal] (0,0) -- (0.8,1.2) -- (1.6,0) -- cycle;
                \fill[brandOrange] (0.8,1.2) -- (1.6,2.4) -- (2.4,1.2) -- cycle;
            \end{tikzpicture}
        };
    \end{tikzpicture}

    \vspace{-0.05cm}
    {\hspace*{1.1cm}\khatt\bfseries\color{brandDark}
        المدرسة العليا للأساتذة ببوزريعة ــ ملحقة بومرداس\par}

    \vspace{1.2cm}
    \begin{minipage}{0.85\textwidth}
        {\color{tealDeep}\khatt\bfseries #1} \\[0.3cm]
        {\raggedleft\color{brandDark}\Huge\khatt\bfseries #2\par} \vspace{0.3cm}
        {\color{brandMuted}\large #3}
    \end{minipage}

    \vspace{1.2cm}
    \begin{minipage}{0.85\textwidth}
        \color{brandMuted}\small
        \begin{tabular}{@{}ll}
            \textbf{الأستاذ:} & د. بن عبد الرزاق يوسف \\
            \textbf{الفئة:}   & طلبة قسم اللغة العربية \\
        \end{tabular}
    \end{minipage}
\end{frame}}

\title{اكتشاف الذكاء الاصطناعي التوليدي}

\begin{document}

\cover{درس الفصل الأول}
{اكتشاف الذكاء الاصطناعي التوليدي}
{مدخل بسيط لطلبة قسم اللغة العربية}

\begin{frame}{أهداف الدرس}
    \begin{tbox}{\faBullseye\quad أهداف الدرس}
        في نهاية هذا الدرس، تستطيع أن:
        \begin{itemize}
            \item تعرّف الذكاء الاصطناعي التوليدي بكلمات بسيطة.
            \item تذكر أنواع المحتوى الذي ينتجه.
            \item تفهم لماذا تتغير الإجابة عندما نغيّر صياغة الطلب.
            \item تتحقق من المعلومة قبل استعمالها في درسك.
        \end{itemize}
    \end{tbox}
\end{frame}

\begin{frame}{ما الذكاء الاصطناعي؟}
    نستعمل كل يوم أجهزة تبدو «ذكية»: الهاتف يقترح الكلمة التالية،
    والخريطة تقترح أقصر طريق.
    وراء هذه الخدمات برامج نسميها \textbf{الذكاء الاصطناعي}.
    \begin{tbox}{\faLightbulb\quad تعريف}
        \textbf{الذكاء الاصطناعي} فرع من علوم الحاسوب،
        يجعل الحاسوب يقوم بمهام تحتاج عادة إلى ذكاء الإنسان،
        مثل الفهم والتعلم والإجابة عن الأسئلة.
    \end{tbox}
\end{frame}

\begin{frame}{ما الذكاء الاصطناعي التوليدي؟}
    كلمة «التوليدي» مأخوذة من الفعل «ولَّد»، أي أنشأ شيئًا جديدًا.
    فالأداة لا تبحث عن إجابة جاهزة فقط، بل \textbf{تنتج محتوى جديدًا}.
    \begin{tbox}{\faLightbulb\quad تعريف}
        \textbf{الذكاء الاصطناعي التوليدي} نوع من الذكاء الاصطناعي
        ينتج نصوصًا أو صورًا أو أصواتًا جديدة، انطلاقًا من طلب يكتبه المستخدم.
    \end{tbox}
    \begin{center}
        \flow[2.3cm]{{المستخدم\\يكتب طلبًا}, {الأداة\\تحلل الطلب}, {محتوى\\جديد}}
    \end{center}
\end{frame}

\begin{frame}{تشبيه يقرّب الفكرة}
    \begin{tbox}{\faChalkboardTeacher\quad التلميذ القارئ}
        تخيّل تلميذًا قرأ آلاف النصوص العربية.
        بعد ذلك، صار قادرًا على كتابة نص جديد بأسلوبه.\\[3pt]
        الأداة التوليدية تشبه هذا التلميذ، لكنها قرأت كمية ضخمة جدًا من النصوص.
    \end{tbox}
\end{frame}

\begin{frame}{ماذا ينتج الذكاء الاصطناعي التوليدي؟}
    \begin{center}
    \begin{tikzpicture}[
        far/.style={draw=brandTeal, fill=brandTeal!12, thick,
            rounded corners=4pt, text width=2cm, align=center,
            minimum height=0.85cm, font=\khatt\bfseries}]
        \node[circle, draw=brandOrange, fill=brandOrange!12, very thick,
            text width=2.3cm, align=center, font=\khatt\bfseries]
            (m) {الذكاء الاصطناعي التوليدي};
        \node[far] (n1) at (0,2.3)     {نصوص};
        \node[far] (n2) at (-4,1.1)    {صور};
        \node[far] (n3) at (-4,-1.1)   {أصوات};
        \node[far] (n4) at (0,-2.3)    {فيديو};
        \node[far] (n5) at (4,-1.1)    {جداول\\ومخططات};
        \node[far] (n6) at (4,1.1)     {برامج\\حاسوبية};
        \foreach \i in {1,...,6} \draw[thick, brandMuted] (m) -- (n\i);
    \end{tikzpicture}
    \end{center}
\end{frame}

\begin{frame}{أمثلة في قسم اللغة العربية}
    \begin{tbox}{\faChalkboardTeacher\quad في قسم اللغة العربية}
        \begin{itemize}
            \item \textbf{نص:} قصة قصيرة لتلاميذ السنة الأولى متوسط.
            \item \textbf{صورة:} رسم يوضح مشهدًا من نص القراءة.
            \item \textbf{صوت:} قراءة مسجلة لقصيدة مع ضبط الحركات.
            \item \textbf{جدول:} جدول يلخص أنواع الخبر مع أمثلة.
        \end{itemize}
    \end{tbox}
\end{frame}

\begin{frame}{أشهر الأدوات}
    توجد أدوات كثيرة، وكلها تقبل الكتابة باللغة العربية:
    \begin{center}
        \begin{tabularx}{0.6\textwidth}{|C|C|}
            \hline
            \ras{الأداة} & \ras{الشركة المطوِّرة} \\
            \hline
            \textenglish{ChatGPT}    & \textenglish{OpenAI} \\ \hline
            \textenglish{Gemini}     & \textenglish{Google} \\ \hline
            \textenglish{Claude}     & \textenglish{Anthropic} \\ \hline
            \textenglish{DeepSeek}   & \textenglish{DeepSeek} \\ \hline
            \textenglish{Qwen}       & \textenglish{Alibaba} \\ \hline
            \textenglish{Perplexity} & \textenglish{Perplexity AI} \\ \hline
        \end{tabularx}
    \end{center}
    الطلب نفسه لا يعطي الإجابة نفسها في كل الأدوات، لذلك من المفيد أن \textbf{نقارن}.
\end{frame}

\begin{frame}{الطلب يغيّر الإجابة}
    نسمي ما نكتبه للأداة \textbf{الطلب} أو \textbf{الأمر} \textenglish{(Prompt)}.
    كلما كان واضحًا ودقيقًا، كانت الإجابة أفضل.

    \vspace{2mm}
    \textbf{\color{brandOrange}طلب عام:}
    \begin{amr}[brandOrange]
        اشرح النعت.
    \end{amr}
    \textbf{\color{tealDeep}طلب دقيق:}
    \begin{amr}[brandTeal]
        اشرح النعت لتلميذ في السنة الأولى متوسط،
        بلغة بسيطة، مع ثلاثة أمثلة من الحياة اليومية.
    \end{amr}
    الطلب الثاني يحدد \textbf{المستوى} و\textbf{اللغة} و\textbf{عدد الأمثلة}،
    لذلك تأتي إجابته أقرب إلى حاجة الأستاذ.
\end{frame}

\begin{frame}{التحقق من إجابات الذكاء الاصطناعي}
    \begin{obox}{\faExclamationTriangle\quad الهلوسة}
        قد تخطئ الأداة وتكتب المعلومة الخاطئة بثقة كبيرة:
        تنسب بيتًا شعريًا إلى غير قائله، أو تعطي إعرابًا خاطئًا.
        \textbf{الإجابة الجميلة ليست بالضرورة إجابة صحيحة.}
    \end{obox}
    نتحقق من المعلومة المهمة بهذه الخطوات:
    \begin{center}
        \flow[1.9cm]{{اقرأ\\الإجابة}, {حدد\\المعلومة}, {ابحث في\\مصدر موثوق}, {قارن}, {قرّر}}
    \end{center}
    \textbf{أمثلة لمصادر موثوقة:} المعاجم المعروفة، كتب النحو المعتمدة،
    الدواوين المحققة، والكتاب المدرسي الرسمي.
\end{frame}

\begin{frame}{كيف يستفيد أستاذ اللغة العربية؟}
    \begin{tbox}{\faChalkboardTeacher\quad استعمالات مفيدة}
        يمكن للأستاذ أن يستعمل الذكاء الاصطناعي التوليدي في:
        \begin{itemize}
            \item اقتراح أمثلة جديدة لقاعدة نحوية.
            \item إعداد أسئلة لفهم نص القراءة.
            \item تبسيط شرح درس في البلاغة.
            \item اقتراح موضوعات للتعبير الكتابي.
        \end{itemize}
    \end{tbox}
    لكن الأستاذ يبقى هو \textbf{المسؤول}: يقرأ، ويصحح، ويختار ما يناسب تلاميذه.
\end{frame}

\begin{frame}{اختبر فهمك}
    \begin{tbox}{\faQuestionCircle\quad صحيح أم خطأ؟}
        \begin{enumerate}
            \item كل ما يقدمه الذكاء الاصطناعي صحيح دائمًا.
            \item يمكن للذكاء الاصطناعي التوليدي إنتاج الصور.
            \item يمكن أن تتغير الإجابة حسب صياغة الطلب.
            \item يجب التحقق من المعلومات المهمة قبل استعمالها.
        \end{enumerate}
    \end{tbox}
    {\small\color{brandMuted} الإجابات: 1) خطأ \quad 2) صحيح \quad 3) صحيح \quad 4) صحيح}
\end{frame}

\begin{frame}{خلاصة الدرس}
    \begin{tbox}{\faCheckCircle\quad خلاصة الدرس}
        \begin{itemize}
            \item الذكاء الاصطناعي التوليدي \textbf{ينتج محتوى جديدًا}:
                  نصوصًا وصورًا وأصواتًا.
            \item جودة الإجابة مرتبطة بجودة \textbf{الطلب}.
            \item الأداة قد تخطئ، لذلك \textbf{نتحقق} دائمًا من مصدر موثوق.
            \item الذكاء الاصطناعي \textbf{مساعد} للأستاذ، وليس بديلًا عنه.
        \end{itemize}
    \end{tbox}
\end{frame}

\end{document}
